%% Dependent Type System
\chapter{Dependent Type System}
\label{ch:dependent-types}

Lambdust's dependent type system extends the gradual typing foundation with types that can depend on values, enabling precise specification of program properties and proof-carrying code through the Curry-Howard correspondence.

\section{Dependent Types Overview}
\label{sec:dependent-overview}

Dependent types allow types to be parameterized by values, not just other types. This enables expressing precise invariants and specifications:

\begin{itemize}
\item \textbf{Pi Types}: Dependent function types $(x : A) \to B(x)$
\item \textbf{Sigma Types}: Dependent pair types $(x : A) \times B(x)$  
\item \textbf{Refinement Types}: Types with logical predicates $\{x : A | P(x)\}$
\item \textbf{Indexed Types}: Types parameterized by values like $\text{Vector}(n, A)$
\item \textbf{Equality Types}: Path equality types $a =_A b$
\end{itemize}

\section{Pi Types (Dependent Functions)}
\label{sec:pi-types}

\indexentry{Pi type}
\indexentry{dependent function}

Pi types generalize function types to allow the result type to depend on the parameter value:

\begin{grammar}
\production{\var{pi-type}}{\produces}{(Pi (\var{parameter} \var{type}) \var{body-type})}{}
\production{}{\alternative}{((\var{parameter} : \var{type}) -> \var{body-type})}{}
\end{grammar}

\begin{example}[Pi Types]
\begin{lambdustcode}
;; Vector lookup with length-dependent result
(define vector-ref :: (Pi (n : Natural) 
                         (-> (Vector n Any) 
                             (i : (Range 0 n)) 
                             Any)))

;; Type-level computation in return type  
(define replicate :: (Pi (n : Natural) (A : Type) 
                         (-> A (Vector n A))))

;; Dependent elimination for natural numbers
(define nat-elim :: (Pi (P : (-> Natural Type))
                        (-> P(Zero)
                            ((n : Natural) -> P(n) -> P(Succ(n)))
                            (n : Natural) -> P(n))))
\end{lambdustcode}
\end{example}

\section{Sigma Types (Dependent Pairs)}
\label{sec:sigma-types}

\indexentry{Sigma type}
\indexentry{dependent pair}

Sigma types generalize pair types where the second component's type depends on the first:

\begin{grammar}
\production{\var{sigma-type}}{\produces}{(Sigma (\var{parameter} \var{type}) \var{body-type})}{}
\production{}{\alternative}{((\var{parameter} : \var{type}) $\times$ \var{body-type})}{}
\end{grammar}

\begin{example}[Sigma Types]
\begin{lambdustcode}
;; Dependent pair with length-indexed vector
(define sized-vector :: (Sigma (n : Natural) (Vector n String)))

;; Create a sized vector
(define my-data :: sized-vector
  (pair 3 #("hello" "world" "!")))

;; Access components
(define get-size :: (-> sized-vector Natural))
(define (get-size sv) (proj1 sv))

(define get-vector :: (Pi (sv : sized-vector) 
                          (Vector (get-size sv) String)))
(define (get-vector sv) (proj2 sv))

;; Existential type (special case of Sigma)
(define abstract-stack :: (Sigma (S : Type)
                                 (& (-> Unit S)           ; empty
                                    (-> S Any S)          ; push  
                                    (-> S (Maybe Any S)))) ; pop
\end{lambdustcode}
\end{example}

\section{Refinement Types}
\label{sec:refinement-types}

\indexentry{refinement type}
\indexentry{predicate}

Refinement types restrict base types with logical predicates:

\begin{grammar}
\production{\var{refinement-type}}{\produces}{(Refinement (\var{variable} \var{base-type}) \var{predicate})}{}
\production{}{\alternative}{\{\var{variable} : \var{base-type} | \var{predicate}\}}{}
\end{grammar}

\begin{example}[Refinement Types]
\begin{lambdustcode}
;; Non-zero numbers
(define NonZero (Refinement (x Number) (not (= x 0))))

;; Safe division
(define safe-divide :: (-> Number NonZero Number))
(define (safe-divide x y) (/ x y))

;; Sorted lists  
(define SortedList 
  (Refinement (lst (List Number)) 
              (sorted? lst)))

;; Insert maintains sortedness
(define sorted-insert :: (-> Number SortedList SortedList))

;; Even numbers
(define Even (Refinement (n Integer) (= (modulo n 2) 0)))

;; Odd numbers  
(define Odd (Refinement (n Integer) (= (modulo n 2) 1)))

;; Adding evens gives even
(define add-even :: (-> Even Even Even))
\end{lambdustcode}
\end{example}

\section{Indexed Types}
\label{sec:indexed-types}

\indexentry{indexed type}
\indexentry{type family}

Indexed types are parameterized by values at the type level:

\begin{example}[Indexed Types]
\begin{lambdustcode}
;; Length-indexed vectors
(define-indexed-type (Vector n A) 
  (vector-data :: (RawVector A))
  (vector-length :: n))

;; Type-safe vector operations
(define make-vector :: (Pi (n : Natural) (A : Type) (-> A (Vector n A))))

(define vector-append :: (Pi (m n : Natural) (A : Type)
                             (-> (Vector m A) (Vector n A) 
                                 (Vector (+ m n) A))))

;; Finite sets indexed by size  
(define-indexed-type (FinSet n)
  (elements :: (Vector n Natural)))

;; Matrix with dimensions
(define-indexed-type (Matrix rows cols A)
  (data :: (Vector (* rows cols) A)))

;; Type-safe matrix multiplication
(define matrix-multiply :: (Pi (m n p : Natural) (A : Type)
                              (=> (Ring A)
                                  (-> (Matrix m n A) (Matrix n p A) 
                                      (Matrix m p A))))
\end{lambdustcode}
\end{example}

\section{Equality Types}
\label{sec:equality-types}

\indexentry{equality type}
\indexentry{path equality}

Equality types represent proofs that two values are equal:

\begin{grammar}
\production{\var{equality-type}}{\produces}{(\var{term1} = \var{term2})}{}
\production{}{\alternative}{(\var{term1} = $_{\var{type}}$ \var{term2})}{}
\end{grammar}

\begin{example}[Equality Types]
\begin{lambdustcode}
;; Reflexivity proof
(define refl :: (Pi (A : Type) (x : A) (x = x)))
(define (refl A x) (equal-refl x))

;; Symmetry proof  
(define sym :: (Pi (A : Type) (x y : A) (-> (x = y) (y = x))))
(define (sym A x y prf) (equal-sym prf))

;; Transitivity proof
(define trans :: (Pi (A : Type) (x y z : A) 
                     (-> (x = y) (y = z) (x = z))))
(define (trans A x y z prf1 prf2) (equal-trans prf1 prf2))

;; Substitution (congruence)  
(define subst :: (Pi (A B : Type) (f : (-> A B)) (x y : A)
                     (-> (x = y) ((f x) = (f y)))))
(define (subst A B f x y prf) (equal-cong f prf))

;; Arithmetic equality proofs
(define plus-zero :: (Pi (n : Natural) ((+ n 0) = n)))
(define plus-succ :: (Pi (n m : Natural) ((+ n (Succ m)) = (Succ (+ n m)))))
\end{lambdustcode}
\end{example}

\section{Universe Hierarchy}
\label{sec:universes}

\indexentry{universe}
\indexentry{type of types}

To avoid Russell's paradox, types are organized in a hierarchy of universes:

\begin{example}[Universe Hierarchy]
\begin{lambdustcode}
;; Universe levels
Type0 :: Type1      ; Base types live in Type0
Type1 :: Type2      ; Type0 itself has type Type1  
Type2 :: Type3      ; And so on...

;; Base types
Number :: Type0
String :: Type0  
Boolean :: Type0

;; Type constructors
List :: (-> Type0 Type0)
Vector :: (-> Natural Type0 Type0)

;; Polymorphic types
(All (A : Type0) (-> A A)) :: Type0     ; Monomorphic polymorphism
(All (A : Type1) (-> A A)) :: Type1     ; Type-polymorphic

;; Universe polymorphism
(define id :: (Pi (l : Level) (A : Type l) (-> A A)))
(define (id l A x) x)
\end{lambdustcode}
\end{example}

\section{Proof Terms and Tactics}
\label{sec:proof-terms}

\indexentry{proof term}
\indexentry{tactic}

Dependent types enable proof-carrying code through the Curry-Howard correspondence:

\section{Basic Proof Construction}
\label{sec:basic-proofs}

\begin{example}[Basic Proofs]
\begin{lambdustcode}
;; Theorem: addition is associative
(define plus-assoc :: (Pi (m n p : Natural) 
                          (((+ (+ m n) p) = (+ m (+ n p))))))

;; Proof by induction on m
(define (plus-assoc m n p)
  (induction m
    ;; Base case: m = 0
    (base-case 
      (trans (+ (+ 0 n) p)          ; Start with goal
             (+ n p)                 ; Simplify (+ 0 n) = n  
             (+ 0 (+ n p))           ; Target form
             (plus-zero n)           ; (+ 0 n) = n
             (sym (plus-zero (+ n p))))) ; (+ 0 (+ n p)) = (+ n p)
    
    ;; Inductive case: m = (Succ k)
    (inductive-step k ih
      (trans (+ (+ (Succ k) n) p)   ; Start
             (+ (Succ (+ k n)) p)    ; Definition of + 
             (Succ (+ (+ k n) p))    ; Definition of +
             (Succ (+ k (+ n p)))    ; Induction hypothesis  
             (+ (Succ k) (+ n p))    ; Definition of +
             (plus-succ-left k n)    ; Lemma
             (plus-succ k (+ n p))   ; Another lemma
             ih))))                  ; Use IH
\end{lambdustcode}
\end{example}

\section{Proof Tactics}
\label{sec:tactics}

High-level tactics automate common proof patterns:

\begin{example}[Proof Tactics]
\begin{lambdustcode}
;; Define a theorem with tactic proof
(theorem plus-comm :: (Pi (m n : Natural) ((+ m n) = (+ n m)))
  (proof-by-induction m
    (base-case 
      (by-simplification))
    (inductive-case k ih
      (by-rewriting 
        (plus-succ-right n k)
        ih
        (plus-succ-left k n)))))

;; Automated tactics
(theorem list-length-app :: (Pi (A : Type) (xs ys : (List A))
                                ((length (append xs ys)) = 
                                 (+ (length xs) (length ys))))
  (auto-induction xs))      ; Automatic induction

;; Interactive proof development
(theorem insertion-sort-correct :: (Pi (xs : (List Natural))
                                      (sorted? (insertion-sort xs)))
  (interactive-proof
    ;; Proof script will be filled in interactively
    ))
\end{lambdustcode}
\end{example}

\section{Program Extraction}
\label{sec:extraction}

\indexentry{program extraction}
\indexentry{Curry-Howard}

Constructive proofs can be automatically extracted to executable programs:

\begin{example}[Program Extraction]
\begin{lambdustcode}
;; Theorem with computational content
(theorem exists-sqrt :: (Pi (x : Positive) 
                           (Exists (y : Positive) ((* y y) = x)))
  (proof
    ;; Constructive proof using Newton's method
    (let ((y (newton-sqrt x)))
      (pair y (newton-sqrt-correct x)))))

;; Extract the computational part
(define sqrt-function (extract exists-sqrt))
;; sqrt-function :: (-> Positive Positive)

;; Theorem: every list can be sorted  
(theorem sort-exists :: (Pi (A : Type) (ord : (Order A)) (xs : (List A))
                           (Exists (ys : (List A)) 
                             (& (sorted? ys) (permutation? xs ys))))
  (proof
    ;; Constructive proof using merge sort
    (let ((ys (merge-sort xs)))
      (pair ys (merge-sort-correct xs)))))

;; Extract sorting algorithm
(define sort-algorithm (extract sort-exists))
;; sort-algorithm :: (Pi (A : Type) (-> (Order A) (List A) (List A)))

;; Verified compiler optimization
(theorem beta-reduce-correct :: 
  (Pi (A B : Type) (f : (-> A B)) (x : A)
      (((lambda (y) (f y)) x) = (f x)))
  (proof (refl (f x))))

(define beta-reduce (extract beta-reduce-correct))
;; beta-reduce :: (Pi (A B : Type) (-> (-> A B) A B))
\end{lambdustcode}
\end{example}

\section{Type-Level Programming}
\label{sec:type-level-programming}

\indexentry{type-level programming}
\indexentry{type family}

Dependent types enable computation at the type level:

\begin{example}[Type-Level Programming]
\begin{lambdustcode}
;; Type-level natural numbers  
(define-type Nat
  Zero :: Nat
  (Succ Nat) :: Nat)

;; Type-level addition
(define-type-function (Plus m n)
  (Plus Zero n) = n
  (Plus (Succ m) n) = (Succ (Plus m n)))

;; Type-level comparison
(define-type-function (Less-Than m n)  
  (Less-Than Zero (Succ n)) = True
  (Less-Than (Succ m) Zero) = False
  (Less-Than (Succ m) (Succ n)) = (Less-Than m n)
  (Less-Than m m) = False)

;; Statically sized data structures
(define-indexed-type (Array n A)
  (data :: (RawArray A))
  (size :: n)
  (constraint :: (= (raw-array-size data) n)))

;; Type-safe array operations
(define array-get :: (Pi (n : Nat) (A : Type) 
                         (-> (Array n A) 
                             (i : Nat) 
                             (constraint :: (Less-Than i n) = True)
                             A))

;; Type-level lists
(define-type (TypeList k)
  TypeNil :: (TypeList k)  
  (TypeCons k (TypeList k)) :: (TypeList k))

;; Heterogeneous lists with type-level schema
(define HList :: (-> (TypeList Type) Type)
(define-type-function HList
  (HList TypeNil) = Unit
  (HList (TypeCons A As)) = (Pair A (HList As)))

;; Example usage
(define schema :: (TypeList Type) 
  (TypeCons String (TypeCons Natural (TypeCons Boolean TypeNil))))

(define person :: (HList schema)  
  (pair "Alice" (pair 30 (pair #t unit))))
\end{lambdustcode}
\end{example}

\section{Totality and Termination}
\label{sec:totality}

\indexentry{totality}
\indexentry{termination}

The dependent type system ensures all functions are total:

\section{Structural Recursion}
\label{sec:structural-recursion}

\begin{example}[Structural Recursion]
\begin{lambdustcode}
;; Structural recursion on natural numbers (always terminates)
(define plus :: (-> Natural Natural Natural))
(define (plus m n)
  (match m
    (Zero n)
    ((Succ k) (Succ (plus k n)))))    ; Recursive call on smaller argument

;; Structural recursion on lists
(define length :: (Pi (A : Type) (-> (List A) Natural)))  
(define (length A lst)
  (match lst
    (() Zero)
    ((cons x xs) (Succ (length A xs)))))   ; Recursive call on tail
\end{lambdustcode}
\end{example}

\section{Well-Founded Recursion}
\label{sec:well-founded}

\begin{example}[Well-Founded Recursion]
\begin{lambdustcode}
;; Division using well-founded recursion on decreasing measure
(define div :: (-> Natural Natural Natural))
(define (div m n)
  (termination-proof (measure m))      ; Prove m decreases
  (if (< m n) 
      Zero
      (Succ (div (- m n) n))))         ; m - n < m when n > 0

;; Ackermann function with lexicographic ordering
(define ack :: (-> Natural Natural Natural))
(define (ack m n)  
  (termination-proof (lex-measure m n))  ; (m,n) decreases lexicographically
  (match m
    (Zero (Succ n))
    ((Succ k) (match n
                (Zero (ack k (Succ Zero)))
                ((Succ j) (ack k (ack (Succ k) j)))))))
\end{lambdustcode}
\end{example}

\section{Coinduction and Corecursion}
\label{sec:coinduction}

For infinite data structures, Lambdust supports coinduction:

\begin{example}[Coinductive Types]
\begin{lambdustcode}
;; Coinductive streams  
(define-coinductive-type (Stream A)
  (stream-head :: A)
  (stream-tail :: (Stream A)))

;; Corecursive stream generation
(define nats :: (Stream Natural)
  (corecursion 0 
    (lambda (n) 
      (stream n (+ n 1)))))

;; Stream map (productive corecursion)
(define stream-map :: (Pi (A B : Type) (-> (-> A B) (Stream A) (Stream B)))
(define (stream-map A B f s)
  (corecursion s
    (lambda (s)
      (stream (f (stream-head s)) 
              (stream-tail s)))))
\end{lambdustcode}
\end{example}

\section{Integration with Gradual Typing}
\label{sec:dependent-gradual}

\indexentry{gradual dependent types}

Dependent types integrate smoothly with the gradual type system:

\begin{example}[Gradual Dependent Types]
\begin{lambdustcode}
;; Gradual refinement from dynamic to dependent
;; Level 1: Dynamic  
(define vector-get 
  (lambda (vec i) (vector-ref vec i)))

;; Level 2: Simple types
(define vector-get :: (-> (Vector Any) Natural Any)
  (lambda (vec i) (vector-ref vec i)))

;; Level 3: Bounded index
(define vector-get :: (-> (vec : (Vector Any)) 
                          (i : Natural)  
                          (constraint :: (< i (vector-length vec)))
                          Any)
  (lambda (vec i) (vector-ref vec i)))

;; Level 4: Full dependent types
(define vector-get :: (Pi (n : Natural) (A : Type) 
                          (-> (Vector n A) (Fin n) A))
  (lambda (n A vec i) (vector-ref vec i)))

;; Automatic contract generation at boundaries
(define untyped-caller
  ;; This automatically gets contract checking
  (vector-get some-vector some-index))
\end{lambdustcode}
\end{example}

\section{Standard Dependent Types}
\label{sec:standard-dependent}

Lambdust provides a standard library of dependent types:

\begin{example}[Standard Dependent Types]
\begin{lambdustcode}
;; Finite sets
(define Fin :: (-> Natural Type))
(define-indexed-type (Fin n)
  (value :: Natural)
  (constraint :: (< value n)))

;; Bounded natural numbers  
(define (Range lower upper) 
  (Refinement (x Natural) (and (<= lower x) (<= x upper))))

;; Non-empty lists
(define (NonEmptyList A)
  (Refinement (xs (List A)) (not (null? xs))))

;; Sorted vectors
(define (SortedVector n A)
  (Refinement (v (Vector n A)) (sorted-vector? v)))

;; Matrices with matching dimensions
(define matrix-multiply :: (Pi (m n p : Natural) 
                              (-> (Matrix m n Number) 
                                  (Matrix n p Number) 
                                  (Matrix m p Number)))

;; Red-black trees with color invariants
(define-indexed-type (RBTree A color height)
  (node-data :: A)  
  (left :: (Maybe (RBTree A black-or-red height)))
  (right :: (Maybe (RBTree A black-or-red height)))
  (color-invariant :: (valid-rb-coloring? color left right))
  (height-invariant :: (valid-rb-height? height left right)))
\end{lambdustcode}
\end{example}

\section{Error Messages and Inference}
\label{sec:dependent-errors}

The dependent type checker provides informative error messages:

\begin{example}[Dependent Type Errors]
\begin{lambdustcode}
;; Index out of bounds
(vector-get #(1 2 3) 5)
;; Error: Index out of bounds
;; Expected: (Fin 3)  
;; Actual: 5
;; Note: Vector has length 3, but index 5 >= 3

;; Refinement violation
(define positive-only :: (-> Positive Natural))
(positive-only -5)
;; Error: Refinement type violation
;; Expected: Positive (x > 0)
;; Actual: -5  
;; Constraint violated: -5 > 0 is false

;; Proof obligation not discharged
(define (unsafe-div :: (-> Natural Natural Natural))
  (lambda (x y) (/ x y)))
;; Error: Missing proof obligation
;; Need to prove: y is not equal 0
;; Suggestion: Use (y : NonZero) instead of (y : Natural)
\end{lambdustcode}
\end{example}

The dependent type system provides powerful tools for program verification while maintaining integration with Lambdust's gradual typing approach, allowing developers to choose their level of specification precision.